\section{3.4 Matrix Operations}

\subsection{Matrix Addition}
Add entry-by-entry. Only defined when A and B have the same size.
If A = [a_{ij}] and B = [b_{ij}] then A+B = [a_{ij} + b_{ij}].

\textbf{Example:}
[1 2; 3 4] + [5 6; 7 8] = [6 8; 10 12]

\textbf{Properties:} A+B = B+A, (A+B)+C = A+(B+C),
A + 0 = A, A + (-A) = 0.

\subsection{Scalar Multiplication}
$cA$ multiplies every entry of A by scalar $c$.

\textbf{Example:} $3 \cdot [1\ 2;\ 4\ 5] = [3\ 6;\ 12\ 15]$

Properties: $c(A+B) = cA+cB$, $(c+d)A = cA+dA$,
$c(dA) = (cd)A$, $1A = A$.

\subsection{Matrix Multiplication}
$(AB)_{ij} = \sum_{k=1}^{n} a_{ik} b_{kj}$

A is $m \times n$, B is $n \times p$, then AB is $m \times p$.
The inner dimensions must match.

\textbf{Example:} A = [1 2; 3 4], B = [0 1; 2 3]
AB: row 1 of A dot cols of B:
$(1)(0)+(2)(2) = 4$, $(1)(1)+(2)(3) = 7$
$(3)(0)+(4)(2) = 8$, $(3)(1)+(4)(3) = 15$
So AB = [4 7; 8 15].

\textbf{Warning:} $AB \neq BA$ in general.
Properties: $(AB)C = A(BC)$, $A(B+C) = AB+AC$,
$AI = IA = A$.

\section{3.5 Matrix Inverses}

\subsection{Inverse of a 2x2 Matrix}
For $A = [a\ b;\ c\ d]$, if $\det(A) = ad - bc \neq 0$:

$A^{-1} = \frac{1}{ad-bc} [d\ -b;\ -c\ a]$

\textbf{Example:} $A = [2\ 1;\ 5\ 3]$
$\det(A) = (2)(3)-(1)(5) = 1$
$A^{-1} = \frac{1}{1}[3\ -1;\ -5\ 2] = [3\ -1;\ -5\ 2]$

Check: $AA^{-1} = [1\ 0;\ 0\ 1]$ (identity).

\subsection{Row Reduction Method (any size)}
Form the augmented matrix $[A | I]$ and row-reduce.
When the left side becomes $I$, the right side is $A^{-1}$.

\textbf{Example:} Find inverse of A = [1 2 0; 2 5 0; 3 7 1]

Set up [A | I] and reduce:
[1 2 0 | 1 0 0]
[2 5 0 | 0 1 0]
[3 7 1 | 0 0 1]

$R_2 \to R_2 - 2R_1$: [0 1 0 | -2 1 0]
$R_3 \to R_3 - 3R_1$: [0 1 1 | -3 0 1]
$R_3 \to R_3 - R_2$: [0 0 1 | -1 -1 1]
$R_1 \to R_1 - 2R_2$: [1 0 0 | 5 -2 0]

So $A^{-1} = [5 -2 0; -2 1 0; -1 -1 1]$.

\subsection{Theorem 7: Invertible Matrix Theorem}
For an $n \times n$ matrix A, these are equivalent:
A is invertible $\Leftrightarrow$ $Ax = 0$ has only trivial solution
$\Leftrightarrow$ RREF of A is $I_n$
$\Leftrightarrow$ A has $n$ pivot positions
$\Leftrightarrow$ $Ax = b$ has exactly one solution for every $b$
$\Leftrightarrow$ columns of A are linearly independent
$\Leftrightarrow$ columns of A span $R^n$
$\Leftrightarrow$ $\det(A) \neq 0$

\section{3.6 Determinants}

\subsection{2x2 Determinant}
$\det([a\ b;\ c\ d]) = ad - bc$

\textbf{Example:} $\det([3\ 2;\ 1\ 4]) = (3)(4)-(2)(1) = 10$

\subsection{Cofactor Expansion}
Cofactor $C_{ij} = (-1)^{i+j} M_{ij}$
where $M_{ij}$ is the minor (det of matrix with row $i$, col $j$ deleted).

Sign pattern: $[+ - +;\ - + -;\ + - +]$

$\det(A) = a_{i1}C_{i1} + a_{i2}C_{i2} + ... + a_{in}C_{in}$
(expand along any row or column; pick the one with most zeros!)

\textbf{Example:} Expand along row 1:
$A = [1\ 2\ 3;\ 0\ 4\ 5;\ 0\ 0\ 6]$
$\det(A) = 1 \cdot \det([4\ 5;\ 0\ 6]) - 2 \cdot \det([0\ 5;\ 0\ 6]) + 3 \cdot \det([0\ 4;\ 0\ 0])$
$= 1(24-0) - 2(0-0) + 3(0-0) = 24$
For triangular matrices: $\det = $ product of diagonal entries!

\textbf{Example 2:} Expand along col 1 (zeros help):
$A = [3\ 0\ 0;\ 1\ 2\ 0;\ 4\ 5\ 6]$
$\det(A) = 3 \cdot \det([2\ 0;\ 5\ 6]) = 3(12) = 36$

\subsection{Theorem 2: Determinants and Invertibility}
$A$ is invertible $\Leftrightarrow$ $\det(A) \neq 0$

Row ops on determinant:
$R_i \leftrightarrow R_j$: det changes sign.
$R_i \to cR_i$: det multiplied by $c$.
$R_i \to R_i + cR_j$: det unchanged.
$\det(AB) = \det(A)\det(B)$
$\det(A^{-1}) = \frac{1}{\det(A)}$
$\det(A^T) = \det(A)$

\section{4.2 Vector Spaces}

\subsection{Vector Space}
A set $V$ with addition and scalar mult satisfying 10 axioms
(closure, commutativity, associativity, identity, inverse, distributive laws).
Examples: $R^n$, $M_{mn}$ (matrices), polynomials, functions.

\subsection{Subspace}
$H \subseteq V$ is a subspace if:
1. $\mathbf{0} \in H$
2. $u, v \in H \Rightarrow u + v \in H$ (closed under addition)
3. $u \in H, c \in R \Rightarrow cu \in H$ (closed under scalar mult)

\textbf{Example:} Is $H = \{(x, y) : y = 3x\}$ a subspace of $R^2$?
1. $(0,0)$: $0 = 3(0)$ Yes.
2. $(x_1, 3x_1) + (x_2, 3x_2) = (x_1+x_2, 3(x_1+x_2))$ Yes.
3. $c(x, 3x) = (cx, 3cx)$ Yes.
So $H$ is a subspace.

\textbf{Example (not a subspace):} $H = \{(x,y): y = x+1\}$
$(0,0)$ not in $H$ (fails test 1). Not a subspace.

\subsection{Solution Subspaces}
The solution set of $Ax = 0$ (homogeneous) is a subspace of $R^n$.
Called the null space of $A$.

\textbf{Example:} $A = [1\ -2;\ 2\ -4]$
Row reduce: $[1\ -2;\ 0\ 0]$, so $x_1 = 2x_2$, $x_2$ free.
Solution: $t[2; 1]$, $t \in R$. This is a line through origin = subspace.

\section{4.3 Span and Linear Independence}

\subsection{Span}
$\text{Span}\{v_1, ..., v_p\} = \{c_1 v_1 + ... + c_p v_p : c_i \in R\}$
This is always a subspace.

\subsection{Linear Independence}
$\{v_1, ..., v_p\}$ is linearly independent if
$c_1 v_1 + ... + c_p v_p = 0 \Rightarrow$ all $c_i = 0$.

Otherwise linearly dependent.

\textbf{Test:} Form matrix with $v_i$ as columns. Row reduce.
If every column has a pivot $\to$ linearly independent.
If any column is free $\to$ linearly dependent.

\textbf{Example:} $v_1 = [1;2;3]$, $v_2 = [2;4;6]$
$v_2 = 2v_1$, so dependent. $2v_1 - v_2 = 0$ (non-trivial relation).

\textbf{Example 2:} $v_1 = [1;0;0]$, $v_2 = [0;1;0]$, $v_3 = [0;0;1]$
Standard basis vectors. Matrix = $I_3$. Pivots in every col. Independent.

\textbf{Key facts:}
Any set with more vectors than entries in each vector is dependent.
Any set containing $\mathbf{0}$ is dependent.
A single nonzero vector is always independent.

\section{4.4 Basis and Dimension}

\subsection{Basis}
A basis for subspace $H$ is a set $\{v_1, ..., v_p\}$ that:
1. Is linearly independent
2. Spans $H$

Standard basis for $R^n$: $\{e_1, e_2, ..., e_n\}$

\textbf{Finding a basis for $R^n$:}
Columns of a matrix form a basis for their span if they are pivot columns.

\textbf{Example:} $A = [1\ 2\ 0;\ 0\ 0\ 1;\ 0\ 0\ 0]$
Pivots in columns 1 and 3.
Basis for Col($A$) = $\{[1;0;0], [0;1;0]\}$.

\textbf{Example 2 (basis for $R^3$):}
Is $\{[1;0;0],[1;1;0],[1;1;1]\}$ a basis?
Form matrix, $\det = 1 \neq 0$, so yes.

\subsection{Dimension}
$\dim(H)$ = number of vectors in any basis for $H$.
$\dim(R^n) = n$, $\dim(\{0\}) = 0$.

\subsection{Basis for Solution Space}
Row reduce $A$. Write free variables parametrically.
Each free variable gives one basis vector.

\textbf{Example:} $A = [1\ 2\ -1\ 0;\ 0\ 0\ 0\ 0]$
Free: $x_2, x_3$. $x_1 = -2x_2 + x_3$.
$x = x_2[-2;1;0] + x_3[1;0;1]$
Basis: $\{[-2;1;0], [1;0;1]\}$, $\dim = 2$.

\section{4.5 Matrix Spaces}

\subsection{Null Space}
$\text{Nul}(A) = \{x : Ax = 0\}$ (solution set of homogeneous system).
$\text{Nul}(A)$ is a subspace of $R^n$ (where $A$ is $m \times n$).
Find by row reducing $[A | 0]$ and solving.

\subsection{Column Space}
$\text{Col}(A) = \text{Span of columns of }A$.
$\text{Col}(A)$ is a subspace of $R^m$.
Basis: pivot columns of $A$ (from ORIGINAL matrix, not RREF).

\subsection{Rank}
$\text{rank}(A) = \dim(\text{Col}(A)) = $ number of pivot columns.

\subsection{Rank-Nullity Theorem}
For $A$ of size $m \times n$:
$\text{rank}(A) + \text{nullity}(A) = n$
where $\text{nullity}(A) = \dim(\text{Nul}(A)) = $ number of free variables.

\textbf{Example:} $A$ is $3 \times 5$, rank$(A) = 2$.
nullity$(A) = 5 - 2 = 3$.
$\dim(\text{Col}(A)) = 2$, $\dim(\text{Nul}(A)) = 3$.

\textbf{Example 2:} $A = [1\ 0\ -2;\ 0\ 1\ 3]$
$n = 3$, rank $= 2$ (2 pivots), nullity $= 1$.
Free: $x_3$. Basis for Nul: $\{[2;-3;1]\}$.
Basis for Col: $\{[1;0],[0;1]\}$ (first 2 cols of $A$).

\section{5.1 Second-Order Linear DEs}

\subsection{General Form}
$ay'' + by' + cy = f(t)$
Homogeneous when $f(t) = 0$: $ay'' + by' + cy = 0$.

\subsection{Theorem 1: Superposition}
If $y_1$ and $y_2$ are solutions of the homogeneous equation, then
$y = c_1 y_1 + c_2 y_2$ is also a solution (for any constants $c_1, c_2$).

\subsection{Theorem 4: General Solution}
If $y_1, y_2$ are linearly independent solutions of $ay''+by'+cy=0$,
then the \textbf{general solution} is $y = c_1 y_1 + c_2 y_2$.

($y_1, y_2$ are linearly independent iff Wronskian $W(y_1,y_2) \neq 0$.)
$W = \det([y_1\ y_2;\ y_1'\ y_2']) = y_1 y_2' - y_2 y_1'$

\subsection{Constant Coefficient: Characteristic Equation}
For $ay'' + by' + cy = 0$, try $y = e^{rt}$:
$\Rightarrow ar^2 + br + c = 0$ (characteristic equation)

\textbf{Case 1: Two real roots} $r_1 \neq r_2$
$y = c_1 e^{r_1 t} + c_2 e^{r_2 t}$

\textbf{Example:} $y'' - 5y' + 6y = 0$
$r^2 - 5r + 6 = 0 \Rightarrow (r-2)(r-3) = 0$
$r_1 = 2, r_2 = 3$
$y = c_1 e^{2t} + c_2 e^{3t}$

\textbf{Case 2: Repeated root} $r_1 = r_2 = r$
$y = c_1 e^{rt} + c_2 t e^{rt} = (c_1 + c_2 t)e^{rt}$

\textbf{Example:} $y'' - 4y' + 4y = 0$
$r^2 - 4r + 4 = (r-2)^2 = 0 \Rightarrow r = 2$
$y = (c_1 + c_2 t)e^{2t}$

\textbf{Case 3: Complex roots} $r = \alpha \pm \beta i$
$y = e^{\alpha t}(c_1 \cos(\beta t) + c_2 \sin(\beta t))$

\textbf{Example:} $y'' + 4y = 0$
$r^2 + 4 = 0 \Rightarrow r = \pm 2i$
$\alpha = 0$, $\beta = 2$
$y = c_1 \cos(2t) + c_2 \sin(2t)$

\textbf{Example 2:} $y'' - 2y' + 5y = 0$
$r^2 - 2r + 5 = 0$
$r = \frac{2 \pm \sqrt{4-20}}{2} = \frac{2 \pm 4i}{2} = 1 \pm 2i$
$y = e^{t}(c_1 \cos(2t) + c_2 \sin(2t))$

\subsection{Initial Value Problems}
Given $y(t_0) = y_0$ and $y'(t_0) = y_1$, plug into general solution
and derivative to get a 2x2 system for $c_1, c_2$.

\textbf{Example:} $y'' - y = 0$, $y(0)=2$, $y'(0)=1$
$r^2 - 1 = 0 \Rightarrow r = \pm 1$
$y = c_1 e^t + c_2 e^{-t}$
$y' = c_1 e^t - c_2 e^{-t}$
$y(0): c_1 + c_2 = 2$
$y'(0): c_1 - c_2 = 1$
$\Rightarrow c_1 = \frac{3}{2}$, $c_2 = \frac{1}{2}$
$y = \frac{3}{2}e^t + \frac{1}{2}e^{-t}$
